% Chapter 7. Update the summary when a research question gets its final answer, and move
% items from future work into the main text when they are done.
\chapter{Conclusion and future work}
\label{ch:conclusion}

\section{Conclusion}

This thesis built an agent workflow for designing graphene quantum dots for \ce{NO2} sensing, in
which an LLM proposes structures and code and quantum chemistry decide which ones are valid and
good. The LLM writes GQDs as hexagon positions on a honeycomb grid. Code converts them into
molecules exactly, checks the closed-shell condition with a matching algorithm, and explains every
rejection with edits that would fix it. With this feedback and the closed-shell rule in the prompt,
99\% of design requests ended in a valid GQD, against 52\% without them (RQ1).

An \ce{NO2} sensing score was defined from GFN2-xTB calculations of adsorption energy and charge
transfer for \ce{NO2} and the competing gases. On 50 dataset GQDs it identified a design target that
was not chosen by hand: a moderately small gap together with physisorption of \ce{NO2}. A DFT check
showed that xTB is reliable for physisorption energies, overbinds chemisorbed \ce{NO2}, inflates
the physisorbed charge and fails for \ce{O2} (RQ2, in part).

In an evolutionary loop with the same xTB budget, both the LLM and random mutation improved on the
best starting GQD: its score was 0.322, and the eight runs reached best scores of 0.331 to 0.356.
Both moved toward GQDs with gaps of about 0.3 to 0.4~eV that still physisorb \ce{NO2} (RQ4). With
three runs per mutator, the LLM's children scored higher on average in every pair of runs, but
random mutation found the best GQD in two of the three pairs, so the LLM did not find better GQDs
than random mutation (RQ3).

Among the best GQDs the score reduces to the xTB charge transfer, and the DFT benchmark of twelve
GQDs showed that this charge is an artefact: in DFT, physisorbed \ce{NO2} takes almost no charge
(0.01~\(e\) against 0.23~\(e\)), and the score loses its ranking (\(\rho = -0.16\)). xTB remains
useful for what it ranks well, the binding of \ce{NO2}, \ce{H2O} and \ce{NH3} and the selectivity
margin (\(\rho = 0.74\) to 0.95), and DFT removes the \ce{O2} concern, since \ce{O2} does not bind
(RQ2). The workflow therefore works as a search method, but the objective it searched needs a new
response term before its best GQDs can be called good sensors. In DFT, the GQDs that keep a
response are those that form a weak, reversible bond to \ce{NO2}.
\pending{The baseline without evolution (E6).}

\section{Future work}
\label{sec:future}

\paragraph{Replace the response term.} The DFT benchmark shows that the xTB charge on physisorbed
\ce{NO2} cannot serve as the response. The candidates are an orbital descriptor such as the
alignment of the GQD's HOMO with the singly occupied orbital of \ce{NO2}, a correction learned from
DFT data, or a new target regime of weak, reversible chemisorption, which in DFT keeps a clear
charge transfer. The selectivity term should use water as the main competitor. The loops should
then be repeated with the new score, and the best candidates re-ranked with DFT.

\paragraph{A sensor model.} A model of the sensing film, with gas coverage at real concentrations
and humidity, adsorption and desorption kinetics, and the resulting change in resistance, would
turn the molecular numbers into a predicted sensor signal and give the search a device-level goal.

\paragraph{Complete the current experiments.} Add the baseline without evolution (E6) and repeat
the comparison of E5b with more runs per mutator.

\paragraph{Check the final candidates more thoroughly.} Probe all symmetry-distinct adsorption
sites and optimise the most promising complexes, including any bonded \ce{O2} complex, at the DFT
level instead of using xTB geometries.

\paragraph{Beyond pure carbon (Stage 3b).} Nitrogen and boron dopants and edge groups such as
\ce{-OH}, \ce{-COOH} and \ce{-NH2} change how strongly gases bind \cite{zhang2009dopants} and may
offer much more headroom than shape alone. This needs a representation that labels individual
atoms, new validity rules (dopants change the electron count and therefore the closed-shell
condition), a new gap predictor or none, and a recalibrated score.

\paragraph{A surrogate for the sensing score.} The logs of the loop contain hundreds of xTB-scored
GQDs with all their descriptors. A model trained on them could replace part of the xTB budget, in
the same way that the gap predictor does for the gap.

\paragraph{Fine-tuning a smaller model.} The logs are designed to become three training sets: edits
that improved the score (supervised fine-tuning), pairs of better and worse edits of the same
parent (preference training with DPO \cite{rafailov2023dpo}), and target-to-structure examples from
the 14{,}674 labelled dataset GQDs and the scored children. A small open model of about 7 to 8
billion parameters, fine-tuned with LoRA \cite{hu2022lora}, could then be tested in the same loop
against random mutation and the general-purpose LLM.

\paragraph{Generative models as further baselines.} Models trained on SELFIES strings
\cite{krenn2020selfies}, graph or diffusion models that work directly on the hexagon grid, and
reinforcement learning in the style of REINVENT \cite{olivecrona2017reinvent} would place the
LLM-guided loop among established methods.
